\section{Dynamic stabilizer codes}\label{sec:dynamic-stabilizer-codes}

Essentially, if we now just add the content of the last section to what we've said so far about static stabilizer codes, we get dynamic stabilizer codes.
A dynamic stabilizer code on $n$ qubits is defined by an \textit{operation sequence} $\mathcal{O}$.
This is a sequence $\mathcal{O} = (\mathcal{O}_0, \mathcal{O}_1, \ldots)$ of sets of operations, each operation being either a Clifford unitary or Pauli measurement.
All the operations within a set $\mathcal{O}_j$ must commute; a Clifford unitary $C$ and a Pauli measurement $M_P$ commute iff $C$ and $P$ commute.
The operation sequence need not be finite.
Dynamic stabilizer codes with finite operation sequences are also called \textit{Floquet codes}.
In this case $\OOO_j$ is still defined for all $j \geq 0$; we just take the subscript to be modulo $|\OOO|$.

This operation sequence $\OOO$ gives rise to sequences $\SSS = (\SSS_{-1}, \SSS_0, \SSS_1, \ldots)$, $\VVV = (\VVV_{-1}, \VVV_0, \VVV_1, \ldots)$ and $\LLL = (\LLL_{-1}, \LLL_0, \LLL_1)$ of stabilizer groups, codespaces and logical Pauli groups respectively, as follows.
We start with $\SSS_{-1} = \{I\}$.
Then $\SSS_{j}$ is the updated stabilizer group that results from performing all the operations in $\mathcal{O}_j$, as per \Cref{subsec:evolution-of-codespace}.
The codespace $\VVV_j$ is always just defined to be the stabilizer subspace $\VVV_{\SSS_j}$ corresponding to $\SSS_j$, and likewise $\LLL_j$ is always just the logical Pauli group $\NNN(\SSS_j)/\SSS_j$.
Note that all these objects are defined for all $j \geq 0$, even if $\OOO$ is finite.

\subsection{Establishment}

In \Cref{subsec:evolution-of-codespace} we said that the rank of the stabilizer group $\SSS_j$ either stays the same or increases by one after applying a Clifford or measuring a Pauli.
Correspondingly, the dimension of the codespace $\VVV_j$ either stays the same or decreases by one, and the new logical Pauli group is either still isomorphic to $\PPP_k$ or is now isomorphic to $\PPP_{k-1}$.
So the sequence of codespace dimensions $(k_{-1}, k_0, k_1, \ldots)$ is a non-increasing sequence of integers, starting at $k_{-1} = n$.
Therefore there is always some integer $T$ such that for all $t \geq T$ the codespace dimension $k_t$ is a constant $k$ -- even if it just means $k = 0$ eventually (though this would then be useless for computing!).
We say the code becomes \textit{established} at time $T$.

\subsection{Codespace maps and logical Pauli group maps}

The operations in a dynamic code induce maps between elements in the sequences of codespaces and logical Pauli groups.
In \Cref{sec:more-stabilizer-formalism} we said that every time we perform an operation $O$ (either applying a Clifford unitary or measuring a Pauli and getting outcome $m$), a linear map $A_m = A_{m, O}$ is applied to our Hilbert space.
We looked at the restriction of this $A_{m, O}$ to the codespace $\VVV_\SSS$.
Now we compose many of these together to get codespace maps
\begin{equation}\label{eq:codespace_maps}
    \mathcal{V}_{-1} \xrightarrow{C_0} \mathcal{V}_{0} \xrightarrow{C_1} \mathcal{V}_{1} \xrightarrow{C_2} \ldots
\end{equation}
where $C_j = \prod_{O \in \OOO_j} A_{m, O}|_{\VVV_{j-1}}$.
If we were being completely explicit we would also parametrise $C_j$ by a vector $\bsym{m}_j$ of outcomes of Pauli measurements in $\OOO_j$, but instead we leave this implicit.

Similarly, in \Cref{sec:more-stabilizer-formalism} we said each such operation $O$ induced a homomorphism $\phi = \phi_O$ on logical Pauli groups.
Hence we get group homomorphisms
\begin{equation}\label{eq:logical_pauli_group_maps}
    \LLL_{-1} \xleftarrow{\phi_0} \LLL_0 \xleftarrow{\phi_1} \LLL_1 \xleftarrow{\phi_2} \ldots
\end{equation}
where each $\phi_j = \prod_{O \in \OOO_j} \phi_O$.

After establishment at time $T$, we never again measure a logical operator of the code (\hyperref[thm:stabilizer_group_update_measurement:case3]{Case 3}).
So for all $t > T$ the codespace map (plus post-measurement renormalisation) is always a unitary map from $\VVV_{t-1}$ to $\VVV_t$, as shown in \Cref{subsec:evolution-of-codespace}.
Likewise, as per \Cref{subsec:evolution_of_logical_pauli_group}, the logical Pauli group homomorphism $\phi_t$ is now always an isomorphism $\LLL_t \to \LLL_{t-1}$ (so if we like we can now think about its inverse $\phi_t^{-1}$ instead, which goes in the `right' direction).

\subsection{Induced logical Cliffords}\label{subsec:induced-logical-cliffords}

One interesting feature of dynamic stabilizer codes is that logical Cliffords can be applied just by performing the sequence of operations in the operation sequence $\OOO$, as discussed in \Cref{subsec:the-induced-homomorphism}.
Of course, there are trivial ways in which this is true -- e.g.\ choose your favourite static stabilizer code with a transversal Clifford gate, then define a dynamic stabilizer code whose operation sequence is `measure the stabilizers of the static code then at some point apply the transversal Clifford gate'.
But there are also very non-trivial ways in which this is true -- e.g.\ in the next chapter we'll meet a dynamically condensed colour code that encodes two logical qubits and performs the logical gate $\pm (H \otimes H)SWAP$ every three timesteps.

These logical Cliffords occur as follows.
Suppose at some post-establishment time $t$ our stabilizer group is $\SSS_t = \langle a_1 S_1, \ldots, a_r S_r \rangle$, then at a later time $t'$ it's $\SSS_{t'} = \langle b_1 S_1, \ldots, b_r S_r \rangle$.
That is, the stabilizer groups at $t$ and $t'$ are \textit{equal up to signs}.
Then the logical Pauli groups $\LLL_t$ and $\LLL_{t'}$ can be presented using the same representatives $\XXX_j$ and $\ZZZ_j$ at both timesteps.
Let $\phi_{t, t'} = (\phi_{t+1} \circ \ldots \circ \phi_{t'})^{-1}$ be the induced isomorphism from $\LLL_{t}$ to $\LLL_{t'}$.
This therefore uniquely (up to global phase) defines a logical Clifford; namely the Clifford $U$ that acts by conjugation on $\ol{\XXX_j}$ and $\ol{\ZZZ_j}$ as $U\ol{\XXX_j} U^\dagger = \phi_{t,t'}(\ol{\XXX_j})$ and $U\ol{\ZZZ_j}
U^\dagger = \phi_{t,t'}(\ol{\ZZZ_j})$.

\subsection{Distance}\label{subsec:dynamic_distance}

We said that a static stabilizer code $\SSS$ has an algebraic distance; the minimum-weight non-trivial logical operator in $\NNN(\SSS) - \SSS$, or equivalently the distance of $\Pi_{\SSS} \circ \Pi_{\SSS}$ under the phenomenological noise model $\FFF$ restricted to the timestep between the two projectors.
The community hasn't settled on a definition of algebraic distance for a dynamic stabilizer code -- see, for example, \Ccite{Fu_2025} which investigates this.
Two generalisations suggest themselves based on our static definition.

The first is the minimum-weight non-trivial logical operator across all timesteps, i.e.\ across all $\LLL_t = \NNN(\SSS_t)/\SSS_t$.
Though immediately this seems to have some problems -- for example, at $t = -1$ we have $\SSS_{-1} = \{I\}$, hence $\LLL_{-1}$ can be presented via single qubit Paulis $X_j$ and $Z_j$, meaning the algebraic distance would always be 1!
So it seems reasonable to restrict this to all $t \geq T$, where $T$ is the establishment time.
But then the next problem is that this only captures logical errors that occur entirely within one timestep -- what if there's a logical error consisting entirely of data qubit errors but supported across multiple timesteps?

The second suggestion fixes the latter problem.
Algebraic distance could be defined as the distance of the circuit $C = \ldots \circ \Pi_{\SSS_t} \circ \ldots \circ \Pi_{\SSS_0}$ under phenomenological noise restricted to acting at certain timesteps.
By the same reasoning as above, we take these restricted timesteps to be $t \geq T$, where $T$ is the establishment time.
Note that the circuit $C$ is infinitely long!
This is necessary -- for a Floquet code (one with finite operation schedule $\OOO$), the pattern of detectors and logical operators eventually repeats, so we could truncate the circuit at some point without losing any essential information.
But for a general dynamic stabilizer code with possibly infinite operation schedule $\OOO$, we can't.

Unlike in the static case, these two definitions are no longer equivalent -- indeed, we said the second definition explicitly fixes a blind spot in the first one.
The second definition is in fact a strict generalisation of the first, in the sense that any non-trivial logical operator in $\LLL_t$ for $t \geq T$ is also a logical error that can occur in the circuit $C$ under $\FFF$.
We take the second suggestion to be our definition of algebraic distance for a dynamic stabilizer code.
But we note that in practice we're not actually that interested in this quantity -- what really counts is the distance of the circuit used to implement any given code.

\subsection{Implementation}

To implement a dynamic stabilizer code, one just takes a system of $n$ qubits and performs all of the operations in $\mathcal{O}_0$, then in $\mathcal{O}_1$, and so on.
Every time we measure some $P$ such that $\pm P$ was already in the stabilizer group, we form a detector.
If instead $P$ anti-commuted with some at least one stabilizer, we move unitarily from some codespace $\VVV$ to a new one $\VVV'$, as per \Cref{prop:A_m_unitary_up_to_normalisation}, and possibly have to update the representatives of our logical operators, as per \Cref{lem:logical_presentation_WLOG_case_2}.
As in the static case, the set of all violated detectors form the syndrome, and this is passed to a decoder, whose job it is to figure out what errors occurred and correct them.

\begin{example}
    Recall \Cref{ex:bacon_shor_stabilizer_group_updates,ex:bacon_shor_codespace_updates,ex:bacon_shor_logical_pauli_group_updates}, in which we repeatedly measured $X_1 X_2$, $X_3 X_4$, $Z_1 Z_3$ and $Z_2 Z_4$, in that order.
    We can now view this as a dynamic stabilizer code.
    Since the only operations will be measurements, we abuse notation slightly and write $P$ rather than $M_P$ to denote the measurement of $P$.
    So the operation sequence is defined as:
    \begin{equation}
        \OOO = (\OOO_1 = \{X_1 X_2, X_3 X_4\}, \OOO_2 = \{Z_1 Z_3, Z_2 Z_4\})
    \end{equation}
    As we've seen in the examples referenced above, this becomes established at time $T=1$, after the first measurement of $Z_2 Z_4$.
    Thereafter, we have:
    \begin{equation}
        \begin{aligned}
            \mathcal{S}_t &= \langle a_t X_1 X_2, b_t X_3 X_4, c_1 d_1 Z_1 Z_3 Z_2 Z_4 \rangle \quad \text{if $t \geq 1$ even} \\
            \mathcal{S}_t &= \langle c_t Z_1 Z_3, d_t Z_2 Z_4, a_0 b_0 X_1 X_2 X_3 X_4 \rangle \quad \text{if $t \geq 1$ odd}
%            \mathcal{S}_t &= \langle m_{t,X_1 X_2} X_1 X_2, m_{t,X_3 X_4} X_3 X_4, m_{1,Z_1 Z_3}m_{1,Z_2 Z_4} Z_1 Z_3 Z_2 Z_4 \rangle \text{if $t \geq 1$ even} \\
%            \mathcal{S}_t &= \langle m_{t,Z_1 Z_3} Z_1 Z_3, m_{t,Z_2 Z_4} Z_2 Z_4, m_{0,X_1 X_2}m_{0,X_3 X_4} X_1 X_2 X_3 X_4 \rangle \text{if $t \geq 1$ odd}
        \end{aligned}
    \end{equation}
    where $a_t$ denotes the outcome of measuring $X_1 X_2$ at time $t$, $b_t$ the outcome of $X_3 X_4$, $c_t$ of $Z_1 Z_3$, and $d_t$ of $Z_2 Z_4$.
    So this code encodes a single logical qubit.
    The codespace maps move us unitarily between the corresponding codespaces:
    \begin{equation}
        \ldots \xrightarrow{C_t} \mathcal{V}_{t} \xrightarrow{C_{t+1}} \mathcal{V}_{t+1} \xrightarrow{C_{t+2}} \ldots
    \end{equation}
    where $C_t = \sqrt{2}\Pi_{a_t X_1 X_2}$ if $t \geq 1$ even, or $\sqrt{2}\Pi_{c_t Z_1 Z_3}$ if $t \geq 1$ odd.
    After establishment, the logical Pauli group can be presented using the same representatives for all $t \geq 1$:
    \begin{equation}
        \LLL_t = \langle
        \overline{i},
        \overline{\XXX} = \overline{X_2 X_4}, \overline{\ZZZ} = \overline{Z_3 Z_4},
        \rangle
    \end{equation}
    The induced isomorphism $\phi_t^{-1}$ for $t > 1$ does nothing particularly interesting; it always maps $\ol{\XXX} \mapsto \ol{\XXX}$ and $\ol{\ZZZ} \mapsto \ol{\ZZZ}$.
    Every two timesteps the stabilizer groups $\SSS_t$ and $\SSS_{t+2}$ are equal up to signs, so we get a well-defined logical Clifford, but this is always just the identity.
    After establishment, every measurement of $X_3 X_4$ and $Z_2 Z_4$ is deterministic, i.e.\ forms a detector.
    Namely, we have detectors $a_{t-2} b_{t-2} a_{t} b_{t}$ for all even $t > 1$ and $c_{t-2} d_{t-2} c_{t} d_{t}$ for all odd $t > 1$.
\end{example}